\documentclass[10pt,a4paper]{amsart}
\usepackage[margin=19mm]{geometry}
\usepackage{amsmath,amssymb,mathtools}
\usepackage[scr=rsfs,cal=euler]{mathalfa}
\usepackage{xfrac}
\usepackage[unicode,hidelinks,bookmarks=false]{hyperref}
\newcommand{\ie}{i.e.\ }
\newcommand{\cf}{cf.\ }
\newcommand{\resp}{respectively\ }
\newcommand{\Subsection}{\subsection}
\providecommand{\nobreakdash}{\nobreak\hbox{-}}
\setlength{\emergencystretch}{2em}
\multlinegap0pt
\usepackage{bm}
\usepackage{bbm}
\usepackage{dsfont}
\usepackage{xspace}
\usepackage{cancel}
\usepackage{mathtools}
\usepackage{amsthm}
\makeatletter
\@ifundefined{theorem}{\newtheorem{theorem}{Theorem}}{}
\@ifundefined{lemma}{\newtheorem{lemma}[theorem]{Lemma}}{}
\makeatother
\makeatletter
\newcommand{\C}{\mathscr {C}}
\newcommand{\HypDim}{{\mathbb{H}^d}}
\newcommand{\Id}{\mathds 1}
\expandafter\ifx\csname assumptionp\endcsname\relax
\newenvironment{assumptionp}[1]{
  \renewcommand\theassumptionalt{#1}
  \assumptionalt
}{\endassumptionalt}
\fi
\newcommand{\conn}[3]{#1 \longleftrightarrow #2\textrm { in } #3}
\newcommand{\dd}{\, \mathrm{d}}
\newcommand{\ela}{\mathbb E_{\lambda}}
\newcommand{\orig}{o}
\newcommand{\pla}{\mathbb P_\lambda}
\newcommand{\tlam}{\tau_\lambda}
\makeatother
\title[Mean-field behaviour of the random connection model on hyper]{Mean-field behaviour of the random connection model on hyperbolic space}
\author{Matthew Dickson, Markus Heydenreich}
\date{Source: arXiv:2505.09025v2 (CC BY 4.0)}
\begin{document}
\maketitle

\noindent\emph{Source and licence.} Mean-field behaviour of the random connection model on hyperbolic space. Version arXiv:2505.09025v2 (CC BY 4.0), distributed under the Creative Commons Attribution 4.0 International licence, as recorded in the arXiv licence metadata for this exact version and shown on the abstract page https://arxiv.org/abs/2505.09025v2. Full rights notice: licenses/arxiv-2505.09025-CC-BY-4.0.txt.\par\smallskip

\noindent\textit{Editorial scope.} Counted unit: the Lemma whose source label is \texttt{lem:nclusterUpper} in the article (source0.tex lines 1827--1834, source label \texttt{lem:nclusterUpper}), delivered together with the article's own complete original proof of it (source0.tex lines 1836--1869, one \texttt{proof} environment of 34 lines). No mathematics is written, repaired or re-derived: statement and proof are the article's own text line for line. Required context retained uncounted: lem:nPointvsTree (lines 1790--1796); eqn:TreeGraphInequality (lines 1820--1823). Every definition, display and label of those lines is retained unchanged. Omitted from this package and not redistributed: 1--1789, 1797--1819, 1824--1826, 1835--1835, 1870--2333. Nothing inside a retained excerpt is omitted or altered: every retained body is delivered exactly as the article's lines are written, so no omission receipt of any kind accompanies this package. Citations: the retained excerpts carry no citation command at all, so no bibliography entry of the article is transcribed and no reference is redistributed. Reference adaptation: none was needed and none was made; every reference inside the delivered unit resolves inside it, so no unresolved reference or citation is produced. Printed slips kept literal: no printed word, symbol, hypothesis or number of the counted statement or of its proof was altered. Layout and class: the article's own class is \texttt{article} with packages including fullpage, amssymb, amsmath, amsfonts, amsthm, paralist, graphicx, color, float, mathtools, tikz, xfrac, inputenc, fontenc; the wrapper is the lane's pinned \texttt{amsart} editorial wrapper at 10pt on A4, which restates the theorem-like environments the retained text needs and adds the article's own macro definitions that the retained text needs (\texttt{\textbackslash C}, \texttt{\textbackslash HypDim}, \texttt{\textbackslash Id}, \texttt{\textbackslash conn}, \texttt{\textbackslash dd}, \texttt{\textbackslash ela}, \texttt{\textbackslash orig}, \texttt{\textbackslash pla}, \texttt{\textbackslash tlam}), with extra wrapper packages (bm, bbm, dsfont, xspace, cancel, mathtools). The delivered numbering is the unit's own. Assets: no retained span contains a figure environment, a graphics-inclusion command, a PDF-page inclusion, a table or a TikZ picture; no image file is copied into this package, the package declares no asset dependency and no image file is redistributed with it. Licence: version arXiv:2505.09025v2 of this article is distributed under the Creative Commons Attribution 4.0 International licence, as recorded in the arXiv licence metadata for this exact version and shown on the abstract page https://arxiv.org/abs/2505.09025v2. A full rights notice is delivered at licenses/arxiv-2505.09025-CC-BY-4.0.txt. Characters: every retained excerpt is pure ASCII, so the delivered engine, XeLaTeX with its default Latin Modern text font, reports no missing character.
\begin{lemma}
\label{lem:nPointvsTree}
    For all $n\geq 2$, $\lambda>0$, and $x_1,\ldots,x_n\in\HypDim$,
    \begin{equation}
        \tau_{n,\lambda}\left(x_1,\ldots,x_n\right) \leq T_{n,\lambda}\left(x_1,\ldots,x_n\right).
    \end{equation}
\end{lemma}

\begin{multline}
    \tau_{n,\lambda}\left(x_1,\ldots,x_n\right)  \leq \sum_{g\in G_n}\lambda^{r(g)}\int_{\left(\HypDim\right)^{r(g)}} \prod_{\left\{z,z'\right\}\in E(g)}\pla\left(\conn{z}{z'}{\xi^{z,z'}}\right) \dd y_1\ldots \dd y_{r(g)}\\
     = T_{n,\lambda}\left(x_1,\ldots,x_n\right)\label{eqn:TreeGraphInequality}
\end{multline}

\begin{lemma}
\label{lem:nclusterUpper}
    For all $n\geq 1$, there exists $C_n <\infty$ such that
    \begin{equation}
        \ela\left[\#\C\left(\orig,\xi^{\orig}\right)^n\right] \leq C_n \chi\left(\lambda\right)^{2n-1}
    \end{equation}
    for all $\lambda\geq0$.
\end{lemma}

\begin{proof}    
    By a binomial expansion and Mecke's formula,
    \begin{align}
        \ela\left[\#\C\left(\orig,\xi^{\orig}\right)^n\right] &= \ela\left[\left(1+\sum_{x\in\eta}\Id_{\left\{\conn{\orig}{x}{\xi^{\orig}}\right\}}\right)^n\right]\nonumber\\
        &= 1+ \sum_{k=1}^n \binom{n}{k} \ela\left[\sum_{\left(x_1,\ldots,x_k\right)\in\eta^{(k)}}\Id_{\left\{\orig\longleftrightarrow x_1 \longleftrightarrow \ldots \longleftrightarrow x_k \text{ in }\xi^\orig\right\}}\right]\nonumber\\
        &= 1  + \sum^n_{k=1}\binom{n}{k}\lambda^k\int_{\left(\HypDim\right)^k}\tau_{k+1,\lambda}\left(\orig,x_1,\ldots,x_k\right)\dd x_1\ldots \dd x_k.\label{eqn:NMomentEquality}
    \end{align}
From Lemma~\ref{lem:nPointvsTree} we therefore have
\begin{equation}
    \ela\left[\#\C\left(\orig,\xi^{\orig}\right)^n\right] \leq 1  + \sum^n_{k=1}\binom{n}{k}\lambda^k\int_{\left(\HypDim\right)^k}T_{k+1,\lambda}\left(\orig,x_1,\ldots,x_k\right)\dd x_1\ldots \dd x_k.
\end{equation}


By translation invariance and `peeling off' leaves from the Greg trees $g\in G_{k+1}$ in definition of $T_{k+1,\lambda}$, we can arrive at
\begin{multline}
    \lambda^k\int_{\left(\HypDim\right)^k}T_{k+1,\lambda}\left(\orig,x_1,\ldots,x_k\right) \dd x_1\ldots \dd x_k  = \sum_{g\in G_{k+1}}\lambda^{k+r(g)}\left(\int_{\HypDim}\tlam\left(\orig,z\right)\dd z\right)^{\#E(g)}\\
    = \sum_{g\in G_{k+1}}\lambda^{\#E(g)}\left(\int_{\HypDim}\tlam\left(\orig,z\right)\dd z\right)^{\#E(g)},
\end{multline}
where we have used the fact that the Greg trees we are summing over have $k+1+r(g)$ vertices in total and (since they are trees) will have $\# E(g)=k+r(g)$. Mecke's formula then leads to
\begin{multline}
    \sum_{g\in G_{k+1}}\lambda^{\#E(g)}\left(\int_{\HypDim}\tlam\left(\orig,z\right)\dd z\right)^{\#E(g)} = \sum_{g\in G_{k+1}}\left(\ela\left[\#\C\left(\orig,\xi^{\orig}\right)\right]-1\right)^{\#E(g)}\\
    \leq\sum_{g\in G_{k+1}}\ela\left[\#\C\left(\orig,\xi^{\orig}\right)\right]^{\#E(g)}.
\end{multline}
Recall that from the iterative construction of Greg trees that for any $k\geq 1$ and $g\in G_{k+1}$, we have $\# E(g) \leq 2k-1$. Since $\ela\left[\#\C\left(\orig,\xi^{\orig}\right)\right]\geq 1$, we therefore have
\begin{equation}
    \lambda^k\int_{\left(\HypDim\right)^k}T_{k+1,\lambda}\left(\orig,x_1,\ldots,x_k\right) \dd x_1\ldots \dd x_k \leq N_{k+1}\ela\left[\#\C\left(\orig,\xi^{\orig}\right)\right]^{2k-1}.
\end{equation}

Combining this inequality with \eqref{eqn:NMomentEquality} and \eqref{eqn:TreeGraphInequality} then gives
\begin{equation}
    \ela\left[\#\C\left(\orig,\xi^{\orig}\right)^n\right] \leq 1 + \sum^n_{k=1}\binom{n}{k}N_{k+1}\ela\left[\#\C\left(\orig,\xi^{\orig}\right)\right]^{2k-1} \leq \left(n+1\right)!N_{n+1}\ela\left[\#\C\left(\orig,\xi^{\orig}\right)\right]^{2n-1},
\end{equation}
as required.
\end{proof}

\end{document}
